\section{Evaluation}
\label{sec:eval}

\subsection{The CronQuestions lane}
\label{sec:eval-lane}

CronQuestions~\cite{saxena2021cronquestions} pairs a temporal knowledge graph
drawn from Wikidata~\cite{vrandecic2014wikidata} with template questions about
it. The lane uses the \texttt{data\_v2} release. Its graph has
\res{cron/kgfacts} facts (subject, relation, object, start year, end year) over
\res{cron/entities} entities; \res{cron/inverted} have a start after their end
and are dropped for every store, because semantica cannot answer any
point-in-time query while one is loaded (\S\ref{sec:eval-found}), which leaves
\res{cron/facts}. The test split has \res{cron/of} questions of five types. Four
are covered: simple entity (\res{cron/n/se}, ``who held $P$ in year $T$''),
simple time (\res{cron/n/st}), before/after (\res{cron/n/ba}) and first/last
(\res{cron/n/fl}), \res{cron/n} in all. Time join (\res{cron/n/tj}) is not:
\res{cron/n/tjtext} of those questions name their second entity only in the
question text, which the annotation does not carry, and the other
\res{cron/n/tjjoin} require an overlap join that this pass did not build.

\paragraph{A plan shared by every store.}
The lane does not parse questions. It turns each question's annotation
(entities, relation, time, before or after, first or last) into calls of three
kinds --- \texttt{at}($e, r$, direction, $T$), what a pair holds in year $T$;
\texttt{history}($e, r$, direction), every interval of a pair; and
\texttt{touch}($e$), every interval into or out of an entity --- and computes
the answer from the returned intervals with one function shared by every store.
The rules were chosen on the validation split, where the control scores
\res{cron/validall}\% over \res{cron/validn} covered questions, and frozen before
the test split ran. The control answers the same calls from dictionaries over
the fact file; its score is the ceiling of the plan. A fact holds from the first
instant of its start year until the first instant after its end year, and a
question about year $T$ asks at \texttt{$T$-07-01T00:00:00Z}.

What the lane measures is therefore a store's answers to temporal lookups under a
fixed plan, and their cost. It is not question answering: no model reads the
question, and the scores are not comparable to those of systems that parse it,
such as the models reported with the dataset.

\paragraph{How the store is driven.}
A test file in the benchmarks repository is compiled into the \texttt{graph}
package by \texttt{go test -overlay}, leaving the cloud tree unmodified. It opens
the store as production does, one connection and no key, and passes every
assertion through \texttt{fromWire}, the admission that \op{POST /v1/graph} calls,
then records batches directly. It skips the schema check, which on this data has
no declaration to check against and would add two reads per distinct relation
per batch.
Each fact becomes two edges with their own \texttt{at} and \texttt{until}: one
from the subject, and its inverse from the object under relation $\sim\!r$,
because \op{resolve} answers one (entity, relation) pair and ``who held $P$''
needs the fact filed about $P$. The evidence is the line of the fact file.
\texttt{at} performs what \op{resolve} performs --- the bounded read, the declared
cardinality, the resolution --- without the HTTP layer and the tenant lookup;
\texttt{history} and \texttt{touch} perform the raw read. Every question is asked
at $\tk$ equal to the clock after the load. Two write clocks are run:

\begin{description}[leftmargin=1.4cm,style=nextline,itemsep=2pt]
\item[wire] every row is written at one server instant, the start of the load,
  as if a caller filed the whole history at once through the API;
\item[replay] each row's write instant is its own \texttt{at}, as if a pipeline
  had filed each fact when it became so. Only code inside the package can set
  this clock.
\end{description}

\paragraph{How semantica is driven.}
Semantica~0.7.0's graph is built by its \texttt{GraphBuilder}, with temporal
support enabled, from relationships carrying \texttt{valid\_from} and
\texttt{valid\_until}. It is queried with two methods of
\texttt{TemporalGraphQuery}: \texttt{query\_at\_time} and
\texttt{query\_time\_range}.
Neither call takes an entity, so an answer depends only on the instant: each of
the \res{cron/sem/distinct} distinct calls the plan needs is made once, and
questions are answered from its output. Latency is taken on the first 50 calls
of each kind, each made in full, and each timed call is checked against the
shared answer; none differed.

\subsection{Results at \texttt{1fffdf408}}
\label{sec:eval-now}

\begin{table}[h]
\centering
\small
\begin{tabular}{@{}lrrrrr@{}}
\toprule
Hits@1, \% & all & simple entity & simple time & before/after & first/last \\
\midrule
\cronrow{cron/kg}{control: the plan over dictionaries}
\cronrow{cron/sem}{semantica 0.7.0}
\cronrow{cron/wire}{\texttt{/v1/graph}, wire clock}
\cronrow{cron/replay}{\texttt{/v1/graph}, replayed clock}
\midrule
\cronrow{cron/first/wire}{first run, wire clock}
\cronrow{cron/first/replay}{first run, replayed clock}
\bottomrule
\end{tabular}
\caption{CronQuestions test split, \res{cron/n} covered questions \meas{}. The
store's rows ran on \res{cron/when}, semantica's on \res{cron/sem/when} and the
first run on \res{cron/first/when} (UTC). The interval on each of the first four
rows' overall score is
\res{cron/kg/allci}, and on before/after \res{cron/kg/baci}. The first run is the
store at \texttt{c47fd3692e9a}, before the change described in
\S\ref{sec:time}; its overall intervals are \res{cron/first/wire/allci} and
\res{cron/first/replay/allci}.}
\label{tab:cron}
\end{table}

All three stores give the control's answer to every covered question: the
number of predictions that differ from the control's is
\res{cron/sem/differ} for semantica, \res{cron/wire/differ} for the wire clock
and \res{cron/replay/differ} for the replayed clock \meas{}. Before/after at
\res{cron/kg/ba}\% is the plan's own ceiling, not a store's error. The wire and
replay rows agree, as Proposition~\ref{prop:filing} predicts for this data: when
a fact became known no longer decides what the store says the world was.

\begin{table}[h]
\centering
\small
\begin{tabular}{@{}lrrrr@{}}
\toprule
& load, s & rows & size, MB & query p50 / p99, ms \\
\midrule
semantica 0.7.0 & \res{cron/sem/load} & \res{cron/facts} & --- & \res{cron/sem/p50} / \res{cron/sem/p99} \\
\texttt{/v1/graph}, wire clock & \res{cron/wire/load} & \res{cron/wire/assertions} & \res{cron/wire/mb} & \res{cron/wire/p50} / \res{cron/wire/p99} \\
\texttt{/v1/graph}, replayed clock & \res{cron/replay/load} & \res{cron/replay/assertions} & \res{cron/replay/mb} & \res{cron/replay/p50} / \res{cron/replay/p99} \\
\bottomrule
\end{tabular}
\caption{Cost on one Apple M1 Max, every store in-process \meas{}. The store's
latency is every one of the \res{cron/wire/timed} calls, from before the read to
the returned rows; semantica's is \res{cron/sem/timed} calls, 50 of each kind.
The machine was shared. One-minute load average at the start, then its peak
through each timed phase --- wire: \res{cron/wire/loadavg}. Replay:
\res{cron/replay/loadavg}. Semantica: \res{cron/sem/loadavg}.}
\label{tab:cost}
\end{table}

\paragraph{Where the cost goes.}
The store's median \texttt{at} call is \res{cron/wire/at/p50}\,ms and its median
\texttt{history} call \res{cron/wire/history/p50}\,ms. Each is one indexed read of
one pair; \texttt{at} bounds the read by the point and resolves the rows it
reads. Semantica's
corresponding medians are \res{cron/sem/at/p50} and \res{cron/sem/history/p50}\,ms.
Its query calls take no entity, and \texttt{query\_at\_time} reconstructs the
whole graph at the instant, beginning with a deep copy, before any filter by
entity can apply. The ratio of the two overall medians is about
\res{cron/ratio}. It is a property of the two interfaces as the lane uses them,
not a comparison of two implementations of the same operation: an index by entity
in front of semantica's reconstruction would change it. Semantica's load is a list
of dictionaries and builds no index, which is why it takes \res{cron/sem/load}\,s;
the store's \res{cron/wire/load}\,s is admission, content addressing, the table's
indexes and the full-text index for \res{cron/wire/assertions} rows.

The load averages in Table~\ref{tab:cost} are high for a ten-core machine and
differ between rows, so every latency is an upper bound for this hardware and
small differences between the two clocks mean nothing. The ratio between the
stores is close to five orders of magnitude and follows from the work each call
does.

\paragraph{Refused rows.}
Under the wire clock, \res{cron/wire/refused} statements were refused because
their start lies after the server clock: Wikidata records terms that have been
planned but have not begun, and admission refuses an \texttt{at} more than five
minutes ahead. No covered question asks about them. Under the replayed clock the
write instant is the statement's own start and nothing is refused, which is the
whole difference between \res{cron/wire/assertions} and
\res{cron/replay/assertions} rows.

\subsection{The first run}
\label{sec:eval-first}

The same lane, run against the store at \texttt{c47fd3692e9a}, is the last two
rows of Table~\ref{tab:cron} \meas{}. That version had one read instant and one
value per pair, so the driver filed each fact as four assertions: the edge and
its inverse at the start, and a retraction of each at the end.

\paragraph{Wire clock: every point-in-time question missed.}
Every row was written at one instant in 2026, so its knowable instant was in
2026. The single \texttt{as\_of} bounded the knowable instant, and an
\texttt{at} call for year $T$ asked at $T$; nothing had been knowable in any year
the questions ask about. The store answered \res{cron/first/wire/se}\% of
simple-entity questions. The other three types were unaffected because they read
history without a bound. The overall score, \res{cron/first/wire/all}\%, is the
share of covered questions that are not simple-entity questions, less the plan's
before/after ceiling.

\paragraph{Replayed clock: one value per pair.}
With each row's write instant set to its own \texttt{at}, transaction time
coincided with valid time and the read at $T$ saw the rows begun by $T$. The
store then answered \res{cron/first/replay/se}\% of simple-entity questions
(interval \res{cron/first/replay/seci}). Table~\ref{tab:misses} gives those
\res{cron/first/replay/misses} misses by relation; the run's other
\res{cron/first/replay/bamiss} misses are the before/after questions the plan
cannot answer. Each of the \res{cron/first/replay/misses} is a pair that holds
more than one value over time: a position with successive holders read from the
position's side, a player's teams, a person's employers. The retraction closing one holder's term was the most recent
statement about the pair and so ended every value it held, including that of
a holder whose term had begun before the retraction and was still running.

\begin{table}[h]
\centering
\small
\begin{tabular}{@{}llr@{}}
\toprule
relation & read from & misses \\
\midrule
position held (P39) & the position & \res{cron/first/replay/miss/position} \\
member of sports team (P54) & the member & \res{cron/first/replay/miss/team} \\
employer (P108) & the employee & \res{cron/first/replay/miss/employer} \\
award received (P166) & the recipient & \res{cron/first/replay/miss/award} \\
spouse (P26) & the person & \res{cron/first/replay/miss/spouse} \\
\midrule
all & & \res{cron/first/replay/misses} \\
\bottomrule
\end{tabular}
\caption{Simple-entity questions the first run missed with the replayed clock,
by the relation and direction of the \texttt{at} call \meas{}. The replay
predictions at \texttt{2e6d394} joined to the plan on the question index.}
\label{tab:misses}
\end{table}

\paragraph{What changed.}
The change at \texttt{1fffdf408} addresses the two causes separately: a second
read coordinate, so that a question about the world is not a question about the
record (\S\ref{sec:time}); and \texttt{until} on the statement with declared
cardinality, so that the end of one term ends only that term
(\S\ref{sec:model}). Filing a term as one statement instead of a statement and
a retraction also halved the rows, from \res{cron/first/wire/assertions} to
\res{cron/wire/assertions}, and the file from \res{cron/first/wire/mb} to
\res{cron/wire/mb}\,MB. The load times of the two runs, \res{cron/first/wire/load}
and \res{cron/wire/load}\,s, were taken at different load averages and are not
compared.

\subsection{Defects found during the correction}
\label{sec:eval-defects}

Runs of the lane against the change, before it was merged, located two
defects, and a test was added for each. The figures of those runs are not in the
repository, so we report the defects without them.

\paragraph{A zero sentinel that is also an instant.}
The first form of the change stored an open \texttt{until} as the integer 0. A
term through 1969 ends at 1970-01-01T00:00:00Z, which is Unix time 0, so every
such term was read back as open and held indefinitely. Open is now SQL
\texttt{NULL}, and a migration adds the column to older files without rewriting
a row. The test \texttt{TestATermEndingIn1969Ends} covers it.

\paragraph{Versions learned together treated as corrections.}
The first form also kept one version per statement, chosen by the order. The
knowledge graph repeats statements --- \res{cron/dupkeys} keys, \res{cron/dupends}
of them with different end years --- and those repeats are written together, so
their versions share a knowable instant and the order was decided by the
content address: one of two terms of the same holder was dropped by the value of
a hash. Versions now correct each other only when one was learned later, and
versions learned at the same instant all stand. The test
\texttt{TestTwoAccountsLearnedTogetherBothStand} covers it, and
\texttt{TestTwoTermsOfOneHolderBothHold} covers the neighbouring case of two
overlapping terms of one holder with different start years, where keeping only
the older term reported the wrong term as the most recently begun.

\subsection{Other findings}
\label{sec:eval-found}

\begin{itemize}[leftmargin=*,itemsep=2pt]
\item \textbf{Year 1 was refused as no instant.} Admission tested whether
  \texttt{at} was Go's zero time, which is 0001-01-01T00:00:00Z. The first run
  refused \res{cron/first/refused/zero} assertions (\res{cron/first/refused/zerofacts}
  facts) for that reason. Absence is now detected at the wire, where it can be
  told apart from the first instant of the calendar.
\item \textbf{A read that stopped did not say so.} \op{GET /v1/graph} returned at
  most 10{,}000 rows with no field reporting that the limit was reached. It now
  reports truncation, as \op{resolve} and \op{neighbors} already did. No read in
  either run reached the limit.
\item \textbf{Semantica accepts an inverted interval and then refuses every
  query.} Its graph builder accepts a relationship whose \texttt{valid\_from}
  follows its \texttt{valid\_until} and reports no rejected relationship; every
  later \texttt{query\_at\_time} on that graph raises \texttt{ValueError}. This is why the lane drops the \res{cron/inverted}
  inverted facts for both stores.
\item \textbf{Semantica overflows on a far end.}
  \texttt{query\_time\_range} with an end of 9999-12-31 raises
  \texttt{OverflowError}, because the day-granularity end is widened past the
  largest representable date. \texttt{TemporalBound.OPEN} works, and the lane
  passes it.
\end{itemize}

\subsection{A second lane: LongMemEval-S}
\label{sec:eval-lme}

LongMemEval~\cite{wu2025longmemeval} gives each of 500 questions its own history
of chat sessions between a user and an assistant, with the sessions that hold
the evidence annotated. The lane in \texttt{brain/longmemeval} of the benchmarks
repository scores session recall over those annotations: ALL@$k$, the share of
questions whose evidence sessions are all in the top $k$, with MRR of the first
evidence session. It compares cosine similarity over all-MiniLM-L6-v2 turn
vectors~\cite{reimers2019sbert}, a session scored by its best turn, with the
retrieval engine in \texttt{brain/} of the same repository, run on the same
cached vectors. The engine's configuration was frozen on the LoCoMo development
split~\cite{locomo}, and nothing was tuned on LongMemEval, so all 500 questions
are reported.

\paragraph{What this lane measures and what it does not.}
It measures that retrieval engine, which combines dense similarity with BM25, a
timeline of session dates, adjacency within a session and a second retrieval
hop. It files no assertion and reads nothing from \texttt{/v1/graph}; the
engine's fact, entity and typed-graph generators propose nothing here, because
LongMemEval has no fact layer. It measures retrieval and not answers, and it uses
the release of the S variant that the dataset card now marks as deprecated in
favour of a cleaned one. We include it because it is the other temporal
benchmark in the same repository, and its per-type rows show where a timeline
helps and where it does not. It is not evidence about the assertion store.

\begin{table}[h]
\centering
\small
\begin{tabular}{@{}lrrrrrrr@{}}
\toprule
& & \multicolumn{2}{c}{ALL@5} & \multicolumn{2}{c}{ALL@10} & \multicolumn{2}{c}{MRR} \\
\cmidrule(lr){3-4}\cmidrule(lr){5-6}\cmidrule(lr){7-8}
type & $n$ & cosine & engine & cosine & engine & cosine & engine \\
\midrule
\lmerow{all}{all}
\midrule
\lmerow{ssu}{single-session user}
\lmerow{ssa}{single-session assistant}
\lmerow{ssp}{single-session preference}
\lmerow{ms}{multi-session}
\lmerow{tr}{temporal reasoning}
\lmerow{ku}{knowledge update}
\midrule
\lmerow{abs}{abstention items}
\bottomrule
\end{tabular}
\caption{LongMemEval-S, session recall over the annotated evidence sessions, run
on \res{lme/when} (UTC) \meas{}. ALL@$k$ is in percent, MRR a fraction. The 30 abstention items are also counted in their types. ALL@5
intervals: cosine \res{lme/all/cosfiveci}, engine \res{lme/all/engfiveci}; ALL@10:
cosine \res{lme/all/costenci}, engine \res{lme/all/engtenci}.}
\label{tab:lme}
\end{table}

Per question, ALL@5 changes on \res{lme/changedfive} of the 500 questions,
\res{lme/gainfive} gained and \res{lme/lostfive} lost, and ALL@10 on
\res{lme/changedten}, \res{lme/gainten} gained and \res{lme/lostten} lost
\meas{}. An exact two-sided sign test on the changed questions, computed here
from those counts, gives $p \approx$ \res{lme/signfive} and \res{lme/signten}.
The gain is largest on multi-session and temporal-reasoning questions. The rank
of the first evidence session falls on \res{lme/firstdown} questions; six of
those are temporal questions phrased in relative time (``four weeks ago'', ``last
Saturday''). The timeline generator fires only on a named month or year, and the
engine does not read the question's date, so it cannot resolve them. The engine
takes \res{lme/p50}\,ms per question at the median and \res{lme/pninetyfive}\,ms
at the 95th percentile.
